\section{Geometric formulations and proof transport}

This part records the objects that carry the proofs, and states their
geometric conclusions separately from their complexity formulations.
The locators identify the relevant framework, rather than attributing
the present rank calculations to the older texts. The supplied cote133
files are transcriptions; no fresh inspection of the manuscript images
is claimed.

\subsection{Framework locators and the precise use made of them}

\paragraph{GD-1: group schemes and character gradings.}
SGA~3, Expos\'e~I, Proposition~4.7.2 identifies representations of
an affine group scheme with comodules. Proposition~4.7.3 gives the
character grading for a diagonalizable group, and the following remark
gives the corresponding graded-algebra formulation for affine schemes
with action. These are the integral objects behind the torus-weight
blocks in BR-001--BR-005. A character grading remains meaningful in
every characteristic. A decomposition into characteristic-zero simple
Schur modules is used only after extension to $\mathbb C$.

\paragraph{GD-2: Schur operations, $K_0$, and the filtration.}
SGA~6, Expos\'e~V, Definition~2.1, Example~2.2(b), and Lemma~2.2.1
give exterior-power operations on the Grothendieck group of locally
projective modules and their behavior on an exact sequence.
Definition~2.4 and equation~(2.4.2) state the product and composition
identities. Expos\'e~VI, section~3.2 gives the special-$\lambda$-ring
assertion for the indicated $K'$ of a quasi-compact scheme.
Expos\'e~V, Proposition~3.2 passes from $\lambda$ to $\gamma$;
section~3.10, equation~(3.10.1) defines the augmentation filtration.
Section~7.1, equation~(7.1.1), and Proposition~7.5 give Adams operations.
These locators are carried with their source-check provenance from
CUE\_GCT's separate dictionary D-S6. The present calculations use
symmetric and Schur operations and exact image modules; they do not
deduce an image rank from a class in the $\gamma$-filtration.

\paragraph{GD-3: projective incidence and coherent ideals.}
EGA~II, Definition~5.4.1, printed p.100, is the properness definition;
Theorem~5.5.3(i), printed p.104, states that a projective morphism is
proper. Sections~2.4--2.5, especially Proposition~2.5.2, give the
associated sheaf construction on $\operatorname{Proj}$.
FGA, Bourbaki Expos\'e~221, Theorems~3.1--3.2 and Proposition~3.8,
place the relevant fixed-Hilbert-polynomial Quot functor in a
projective parameter space and its Grassmannian realization.
These exact locators and their inspection scope are carried from
CUE\_GCT D-Proper and D-FGA. Here the incidence varieties are explicitly
closed subvarieties of finite products of Grassmannians, so properness
applies directly to their projections. No effective bound for equations
is inferred from representability alone.

\paragraph{GD-4: derived restriction where it is actually used.}
SGA~6, Expos\'e~III, Definition~1.5 and Notation~1.6 give Tor
independence and the derived tensor product. Expos\'e~I,
Proposition~5.1, Definition~5.2 and paragraph~5.4 give the flat-amplitude
and resolution-length language. These are the locators for the specific
graded Tor calculation in QC-04. A vanishing in one indicated degree
is not a statement of Tor independence in every degree.

\paragraph{Archive dictionary.}
The supplied cote133, scan028, lines6--17 concerns tensor powers and
symmetric operations, and lines19--24 raises the direct-summand
question. Scan019, lines45--56 distinguishes equivariant information
from its rank or augmentation. Scan006, lines25--26 is the universal
product-operation context. The actual character computations below
retain the full Schur component; the subsequent restriction maps still
have to be computed. This is precisely the distinction between a
representation class and an equivariant coherent ideal.

\subsection{Support and multiplicity findings}

\paragraph{S-001: a flag-minor support statement for Schur sections.}
On a rank-$r$ vector bundle with a chosen complete flag, the Cauchy
construction realizes a highest-weight section by products of flag
minors. A coefficient coordinate whose degree outside the first
$j-1$ flag positions exceeds $\lambda_j$ cannot occur in such a
section. The proof applies to the full highest-weight multiplicity
space. It takes place in the polynomial representation category over
characteristic zero, or in the corresponding Schur construction after
the relevant symmetrizer denominators are inverted. GD-1, GD-2 and
cote133 scan028 are the dictionary; the combinatorial minor proof is
the additional mathematical content.

\paragraph{S-002--S-003: the LMR equivariant kernel and image.}
Let $S$ be the homogeneous coordinate algebra of cubic forms and let
$I_D,I_P$ be the ideals of the determinant and permanent images.
In degree12 and type $(19,7,2^5)$ the highest-weight restriction map
to $S/I_D$ has a one-dimensional kernel, explicitly spanned by
$2P_0+21P_1+7P_2+35P_3+7P_4$, and a five-dimensional image.
The corresponding map to $S/I_P$ has zero kernel on the six-dimensional
ambient multiplicity space. Sheafifying puts the equation in the
equivariant ideal sheaf on the projective coefficient space. This is
an assertion about actual kernels and sections in GD-3, with the Schur
grading of GD-1--GD-2. The finite minors prove the stated dimensions;
an equality in $K_0$ alone would not prove them.

\paragraph{S-004: support under an independent padding line.}
The padding construction uses a line bundle summand independent of
the rank-$m^2$ bundle of permanent variables. Its essential bundle
rank is therefore $m^2+1$. The row-cap and nested-support arguments
must be applied to that bundle. S-004 supplies the corrected integer
inequalities for the extra possible row. The change concerns the
underlying coefficient morphism and Schur support, not a change in
the universal $\lambda$-ring identities. Its framework is GD-1--GD-2
and the projective multiplication morphism in GD-3.

\paragraph{S-005--S-006: localization followed by a complete image calculation.}
On the principal open set $D(F_0)$ of the nonic coefficient scheme,
centering is an explicit quotient by a split unipotent group.
Powers of $F_0$ clear the denominators of centered tensor contractions,
producing homogeneous sections of the required Schur type.
The 804-function matrix computes the entire type $(192,8,8,8)$
restriction image on the determinant scheme and shows its kernel is
zero. Injectivity of multiplication by $F_0$ on an integral coordinate
ring propagates that conclusion to every degree in
$(9d-24,8,8,8)$. The hypersurface equation of the octic-polar projection
therefore has tail character at least9. This is a local coordinate and
coherent-ideal calculation in GD-3, followed by the Schur decomposition
of GD-1--GD-2. It does not identify the full defining ideal of the
determinantal locus.

\subsection{Differential certificates and explicit sections}

\paragraph{QC-01--QC-03: rank strata of the permanent morphism.}
The permanent-pencil parameter map is defined over $\mathbb Z$.
Its differential is a map of finite free modules on that affine
parameter scheme. Each certified nonzero minor specifies a principal
open rank locus; its reduction modulo101 or1021 proves the corresponding
minor is nonzero over $\mathbb Q$. The insertion construction propagates
the finite endpoint through the stated family of parameter morphisms.
The torus action supplies the matching fiber-dimension bound.
The exact generic dimensions follow from these concrete tangent maps
and fiber calculations, in GD-1 and GD-3. They are not consequences of
purity or a random-point heuristic.

\paragraph{QC-04: the specified syzygy comparison.}
The multiplication map and the invoked linear-syzygy resolution give
the asserted homogeneous piece of $\operatorname{Tor}_1$. Its geometric
meaning is the stated obstruction group for derived restriction, in
GD-4. Neither the computation nor its dictionary supplies a vanishing
of all higher Tor sheaves.

\paragraph{QC-005--QC-006: an explicit section outside both ideals.}
The thirteen-bracket expression is an integral Schur-functor section
of degree6. It is nonzero at a point of the linear-factor locus and at
an explicit block determinant. Because the ambient multiplicity of
its type is one, both restriction images have multiplicity one.
Thus this line belongs to neither equivariant ideal. The integer
evaluations, rather than only its $K_0$ class, decide this particular
candidate. The relevant dictionary is GD-1--GD-3 and cote133 scan028.

\paragraph{QC-007: arithmetic points for an equivariant restriction map.}
The traceless integer pencils are sections of the determinant parameter
scheme with $F_0=1$ and $F_1=0$. Exact Newton interpolation computes
their normalized tensor coordinates. These points evaluate the actual
sections used in S-006; their nonzero minor establishes injectivity
after scalar extension to characteristic zero. They have no separate
claim to an $\ell$-adic trace or a motivic interpretation.

\paragraph{OBS-001--OBS-004: coefficient projections and all-copy injections.}
Nested coefficient projections are morphisms of affine coefficient
schemes adapted to a complete flag. Their checked Jacobian minors
make the indicated projections dominant. A dominant projection induces
an injection of its polynomial coordinate algebra into the corresponding
image coordinate ring. S-001 places the entire highest-weight space
in that algebra, yielding the claimed full multiplicity equalities.
The plethysm sieves in OBS-002--OBS-003 compute the relevant Schur
classes through symmetric operations and Weyl extraction; the projected
rank certificates provide the additional ideal information. This
separates GD-2 from GD-3 exactly as required by cote133 scan019.

\subsection{Border-rank findings over the arithmetic base}

\paragraph{BR-001--BR-002: proper incidence in every fiber.}
The spaces that may occur as limits of spans of pure tensors are
parameterized by closed incidence schemes in products of relative
Grassmannians. Tautological bundles and exterior-power maps define
the wedge and Koszul degeneracy conditions over the arithmetic base.
The torus decomposes the explicit quotient modules into character
blocks, using GD-1. Proper projection carries limits to incident
subspaces, using GD-3. The proof checks the required fibers over
$\mathbb Z[1/6]$, $\mathbb F_2$ and $\mathbb F_3$ separately.
This establishes the uniform CP-border-rank bound for the specified
integral coefficient tensors, without an integral semisimple Schur
decomposition or an assumption that arbitrary schematic images commute
with every base change.

\paragraph{BR-003: an actual nonunit specialization.}
The displayed determinant2 minor has rank4 over $\mathbb Q$ and rank3
over $\mathbb F_2$. It locates a genuine change of rank in an integral
bundle map, and disproves the proposed literal unit replacement of
that full block. The separate fiber argument of BR-002 is what retains
the global lower bound. A representation-ring identity does not remove
this arithmetic rank drop.

\paragraph{BR-004--BR-005: flag incidences and exact exceptional fibers.}
Invariant complete flags refine the proper incidence of a limiting
seventeen-plane. The inclusion of each partial tautological bundle
gives an additional tensor whose exterior flattenings must satisfy
the prescribed rank bounds. A first nonrigid flag choice is a
projective line in a common root-operator kernel. Pulling a
determinantal rank locus back to that line gives explicitly factored
minor polynomials. Their generic complements and every exceptional
fiber, including infinity and quadratic residue fields over
$\mathbb Q$, are handled separately. This proves the stated finite
exclusions in characteristic zero. The same GD-1/GD-3 setting carries
the proof; no flatness, smoothability or membership in the border
apolarity component is inferred from merely passing necessary tests.

\subsection{Further polynomial kernels and their geometric meanings}

\paragraph{S-007--S-008 and QC-010: injectivity on a multiplication image.}
The projective multiplication morphism sends a pair $(L,Q)$ to
$L^{n-m}Q$; its affine cone gives the actual padded-form image.
On the principal open chart $F_0\ne0$, the normalized coefficient
change is a quotient by the translation subgroup of the chosen
parabolic. The resulting explicit pullback sends the last centered
tensor to a polynomial in a line and freely chosen lower tensors.
The quadratic polarization and cubic-symbol arguments prove that
its kernel contains no section of the stated Schur types. This is
injectivity on the actual homogeneous coordinate algebra, not on
the normalization alone. The inclusion from the degree16 permanent
block specialization gives the same kernel conclusion for the
specified types at $(17,23)$. Its dictionary is GD-1--GD-3 and
cote133 scan028, with the distinction between rank and class in scan019.

\paragraph{S-009: a nonzero equivariant ideal section.}
For binary cubics, the discriminant times the square of the first
coefficient is an explicit degree-six section of the ideal of the
multiple-root multiplication image. In type $(12,6)$ the restriction
map has a one-dimensional kernel and a one-dimensional image inside
a two-dimensional ambient multiplicity space. This concrete coherent
ideal calculation refutes the stated sufficient branch of the inspected
arXiv version. No statement about a normalization or about a numerical
$K_0$ class removes this kernel. The proof transports by GD-1--GD-3;
its version and dependency scope are part of the finding.

\paragraph{QC-008--QC-009: exact sections on arithmetic charts.}
The integral bracket sections and integer parameter points define
matrices for genuine restriction morphisms. Lane arithmetic is only
an exact encoding of their residues; scalar minors determine their
rank. The quadratic-normalized sample chart is defined over
$\mathbb F_{101}$ and is lifted to genuine integer source points.
No lifting of its quadratic-normalization condition to characteristic
zero is asserted. Nonzero minors lift because the coefficient
normalization denominators are units at101. This is a rank-stratum
calculation for maps of locally free modules in GD-1--GD-3.

\paragraph{OBS-005--OBS-006: saturation and a computed padded kernel.}
Multiplication by the nonzero section $F_0$ is injective on an integral
coordinate algebra. The explicit localization identifies a stabilized
Schur multiplicity space and transports its zero determinant kernel
to all admissible degrees in the specified fixed-tail families.
The fourth-direction calculation similarly proves a zero kernel for
the padded pullback, by injectivity of a two-dimensional multiplicity
map and the remaining polarization maps. It does not compute the
determinant ideal in that type. Both are statements about coherent
ideals and their graded Hilbert functions, in GD-2--GD-3.

\paragraph{BR-006--BR-007: specialization and an explicit arc.}
BR-006 proves emptiness of the remaining
relative flag incidence over rational parameter curves. Gaussian
localization specifies a principal open subset; every removed closed
fiber and infinity are recomputed. The consolidated certificate
closes every branch. BR-007 is an
explicit Laurent arc in the space of nineteen pure tensors. Its sum
has a regular limit equal to the ten-plane tautological tensor.
The definition of the secant variety as a Zariski closure puts that
limit in the nineteenth secant variety. These are GD-1/GD-3 arguments,
with no appeal to purity or to an unproved flatness assertion.

\paragraph{QC-011: vanishing of an initial coherent-ideal segment.}
The multiplication morphism $E\times\operatorname{Sym}^{n-1}E\to
\operatorname{Sym}^nE$ has an image whose homogeneous ideal is zero
in every degree at most$n$. Polarization identifies an ideal section
with a multilinear form, and successive injective weight-shift
operators show that this section must vanish identically. This is
an actual graded-kernel statement in GD-3, with the representation
and Schur-functor interpretation in GD-1--GD-2 and cote133 scan028.
It supplies no projective-normality assertion.

\paragraph{OBS-007: a relative second fundamental form.}
On the pivot open of the determinant parameter map the differential
has a locally free left kernel. Its quadratic coefficient on
source directions annihilating the lower projection defines the
relative second fundamental form. Full span of its eighty normal
directions forces every ideal section affine in the last coefficient
tensor to have zero gradient on a dense open set. Dominance of the
lower projection then makes every coefficient of that section zero.
The Schur decomposition identifies the entire indicated highest-weight
space with such affine sections. This transports the proof by
GD-1--GD-3, using explicit localizations and locally free kernels.
The modular kernel vectors lift by pivot solving; arbitrary kernel
base change is not assumed. The conclusion is equality of actual
graded Hilbert multiplicities in every outer degree.

\paragraph{BR-008--BR-009: a permanent flag incidence and finite torus cover.}
The verified solvable stabilizer of the distinguished permanent
quadratic tensor acts on the closed flag-incidence scheme of possible
limiting spans. The complete proper-incidence exclusions prove its
rank18, while an explicit finite-difference arc supplies the upper
bound. For the global question the row-column torus has distinct
weights on the36 quadratic slices, so a nonempty closed low-rank
locus on the Grassmannian has a coordinate fixed point. Exact finite
group orbits and Koszul rank loci reduce it to595 remaining points
up to symmetry. The cover is exhaustive; passing a necessary test
does not assert membership in a secant variety. These are the
GD-1/GD-3 group-scheme, tautological-bundle and proper-incidence
formulations, with the arithmetic scope stated in each finding.

\paragraph{S-010: a complete Schur restriction and hypersurface character bound.}
The2386 integral graph sections form the entire indicated stabilized
Schur multiplicity space. Their evaluation at genuine integer pencil
points has a nonzero maximal minor, so the equivariant coherent ideal
has zero intersection with this space. Multiplication by the nonzero
section $F_0$ and the complete centered dimension transport the zero
kernel to every admissible outer degree. The same zero kernel excludes
a character-nine generator of the first-eight-jet hypersurface ideal.
This is GD-1--GD-3 and cote133 scans028 and019, using actual sections
and image ranks, with no replacement by a numerical representation class.

\paragraph{BR-010: refining a proper incidence by closed linear strata.}
The explicit connected solvable subgroup acts on the same proper
flag incidence. Every retained unipotent root is checked on the actual
tensor. All-coordinate-endpoint substitution computes a Grassmannian
minimum using torus fixed points; an individual strong endpoint would
not suffice. Intersections of first-wedge kernel subspaces give a
finite exact cover of rank strata. The conservative union removes only
complete exclusions and retains every unfinished incidence. Its
geometric setting remains GD-1 and GD-3, with characteristic-zero
localizations and exact closed exceptional fibers.

\paragraph{OBS-008--OBS-009: differential closure of Schur sections.}
The high-coordinate change is an affine isomorphism over the lower
coefficient chart. Equivariant comultiplication and Littlewood--Richardson
containment show that the bounded first-row space is closed under all
high partial derivatives and that its high Hessians lie in the explicitly
listed quadratic Schur bundle. Independent normal and quadratic Taylor
forms force each ideal section to have zero gradient on a dense open.
Induction on its high degree then gives a zero actual ideal kernel.
This is a coherent-ideal and locally free bundle calculation in
GD-1--GD-3, with the Schur operations of cote133 scan028 and the
class-versus-image distinction of scan019. It proves the stated
all-degree Hilbert multiplicity equalities and the hypersurface
character bound; no numerical class is substituted for a restriction map.

\paragraph{QC-012: an explicit complete three-dimensional image.}
Three integral bracket sections of the same Schur type evaluate to a
matrix with nonzero determinant46 modulo101 at three actual integer
block pencils. The ambient multiplicity is independently three, so
this is a complete equivariant image computation. QC-011 supplies
zero padded kernel in the same degree. The result concerns the
actual graded ideals in GD-3 and the group-scheme/Schur constructions
of GD-1--GD-2 and cote133 scan028.

\paragraph{S-011: coefficient functionals on a closed image.}
The block determinant construction maps the ternary septic coefficient
scheme into the quaternary nonic determinant closure. Coefficients in
auxiliary scalar parameters therefore define linear functionals on its
actual coordinate algebra. Three such functionals and one integer pencil
evaluation have a nonzero four-square minor on four integral Schur
sections. The ambient character calculation proves that these sections
exhaust the multiplicity space. Thus the corresponding equivariant
coherent ideal has zero kernel in this degree and type. This is
GD-1--GD-3, with the full-image distinction of cote133 scan019 and
the Schur construction of scan028; no normalization replaces the image.

\paragraph{OBS-010: a larger locally free normal and Hessian bundle.}
The relative Taylor construction of OBS-008 extends to three high
coefficient tensors. Its normal bundle has rank480 and its allowed
quadratic Schur bundle has rank2740. A nonzero3220-square minor of
their combined bilinear map specifies a principal open rank locus.
Differential closure then gives a zero equivariant ideal intersection
in every indicated homogeneous degree. The dictionary is GD-1--GD-3
and cote133 scans028 and019. Only the asserted arithmetic pivot open
is used; no general base-change theorem for image kernels is assumed.

\paragraph{BR-011: exact closed covers of projective rank loci.}
The parameter spaces for nonrigid invariant flag steps are projective
linear spaces. Explicit homogeneous minors cut out necessary rank loci.
Factoring these polynomials gives a cover by closed hypersurfaces;
linear components are recursed into, and every unsupported nonlinear
component is retained. This is a conservative calculation in the
proper incidence of GD-3 with the group-scheme weight decomposition
of GD-1. The generic numerical sample only selects a minor; its exact
polynomial and complete closed exceptional cover carry the proof.

\paragraph{BR-012: a simultaneous proper incidence and quotient weights.}
The closure of the rank-one span incidence in a product of four
Grassmannians supplies compatible pairwise and triple limiting
tautological subspaces. A fixed point for the connected solvable
stabilizer makes them simultaneously invariant. After the two unique
pairwise extensions are fixed, their intersection defines a
27-dimensional vector space and its contractions define a
36-dimensional quotient ambient space. Weightwise ranks of the
projection maps bound every possible remaining subspace, for all
dimensions ten through seventeen. The rational inequality
$4\dim W\le65$ contradicts the required rank17 tautological
subspace. This is GD-1/GD-3 and cote133 scan019, using actual
intersections and ranks, not just classes in $K_0$.

\subsection{Frameworks that supply no additional step here}
The finite-field numbers in this packet are residues of explicit
integer or rational coefficient calculations. They are not point counts
or Frobenius traces. Consequently SGA~4, SGA~5 and Weil~II weight or
purity theorems are not inputs to these proofs, and no motive is
constructed. The relevant transport consists of group schemes,
Schur operations, coherent ideal sheaves, localization and proper
incidence. Introducing cohomological weight language would not decide
an uncomputed orbit-coefficient rank.
